%=======================================================================
%  CHAPTER 5 — STRUCTURED KNOWLEDGE REPRESENTATION : SOLVED PROBLEMS
%=======================================================================
\section{Structured Knowledge Representation \textcolor{sub}{\large -- Solved Problems}}

{\color{sub}\footnotesize Every problem opens with the question
\mk{exactly as the paper prints it}, the paper's own spelling and
slips included. \gap$\cdot$\gap \textbf{11 sentence sets over 12 papers}, plus the
net-to-frame conversion and the net/frame/FOPL question. \gap$\cdot$\gap Where a lecture
deck or the textbook already draws the answer, the published figure is printed and the
caption says exactly what differs from the paper's wording. Where none does, the answer is
given as an \textbf{arc list} $-$ every node, every link, every direction $-$ from which
the drawing is mechanical.}

\lead{\mk{The four-step recipe.} It works on all eleven sets:}
\begin{itemize}
  \item \textbf{Step 1. Underline every noun phrase.} Each distinct one is a
        \textbf{node}. Do not invent nodes and do not merge two that the paper keeps
        apart.
  \item \textbf{Step 2. Sort the nodes into classes and individuals.} A proper name, or
        anything the sentences call \emph{a} something, is an \textbf{individual}. Anything
        that other things can \emph{be} is a \textbf{class}.
  \item \textbf{Step 3. Draw the class spine first, bottom to top}, joined by \code{isa}
        arcs, then hang each individual off its class with an \code{instance} arc.
  \begin{itemize}
    \item \mk{Class to class is \code{isa}. Individual to class is
          \code{instance}.} Getting these two the wrong way round is the single
          commonest lost mark in this question.
  \end{itemize}
  \item \textbf{Step 4. Add one labelled arc per remaining sentence}, always
        \emph{out of} the thing the sentence is about. The verb becomes the arc label.
  \item \textbf{Then check three things}, which is what separates a 6 from a 4:
  \begin{itemize}
    \item \textbf{Count.} One arc per sentence, plus the spine. If a sentence produced no
          arc you have lost a mark.
    \item \textbf{Push properties up.} If two individuals of the same class both carry a
          property, draw it \textbf{once on the class} and let them inherit. Say in one
          line that you did, and why.
    \item \textbf{Flag what the net cannot say.} Quantifiers, negation and disjunction do
          not draw. A sentence like \emph{``every mammal gives birth to a baby''} gets an
          arc out of the class node and \mk{a note saying the arc is read
          as $\forall x \in Mammal$}. Examiners give marks for that note.
  \end{itemize}
  \item \textbf{Layout that always works}: most general class at the \textbf{top},
        \code{isa} arcs pointing \textbf{up}, individuals on a row near the bottom, and
        their own properties fanning \textbf{sideways}. Every published figure in this
        chapter is drawn that way.
\end{itemize}

\hr

%=======================================================================
\T{5.N.1 A person is a mammal: Sakti Gauchan, and the same set as Sandeep Lamichhane}

\creamq{Convert the given sentences into a semantic net}{\m{7}\m{2+6}}
{\color{sub}\footnotesize \tP{2} \yr{72 Ma \m{7}, 78 Po \m{2+6}} \gap$\cdot$\gap
two papers, \mk{one net}: \yr{78 Po} is \yr{72 Ma} with the cricketer's
name changed \gap$\cdot$\gap both also ask \emph{``what are frames and semantic net''}
first, which is \S5.6 of the theory}

\begin{asked}{72 Ma \textperiodcentered{} Q6 \textperiodcentered{} \m{7}}
What are Frames and Semantic Net? Convert the given sentences in Semantic Net.
\begin{itemize}[topsep=1pt]
  \item i) A person is a mammal.
  \item ii) Sakti Gauchan is a person.
  \item iii) Person has nose.
  \item iv) Sakti Gauchan is in Nepalese team.
  \item v) Uniform color of sakti Gauchan is Red/Blue.
\end{itemize}
\end{asked}

\begin{asked}{78 Po \textperiodcentered{} Q6 \textperiodcentered{} \m{2+6}}
What are Frames and Semantic Net? Convert the given sentences in Semantic Net.
\begin{itemize}[topsep=1pt]
  \item A person is a mammal.
  \item Sandeep Lamichane is a person.
  \item Person has nose.
  \item Sandeep Lamichane is in Nepalese cricket team.
  \item \mk{Uniform color of Sakti Gauchan is Red/Blue.}
\end{itemize}
{\color{sub} The fifth sentence names a \emph{different} cricketer: the paper was
copied from the \yr{72 Ma} set and this one line was never updated. Reproduced as
printed.}
\end{asked}

\sbsr{0.53}{%
\lead{Step 1--2: nodes}
\begin{itemize}
  \item \textbf{Classes}: Mammal, Person, Nepalese team (a named team, so it can also be
        drawn as an individual $-$ either is accepted).
  \item \textbf{Individual}: Sakti Gauchan.
  \item \textbf{Values}: Nose, Red/Blue.
\end{itemize}
\lead{Step 3--4: the arcs}
\begin{tabularx}{\linewidth}{@{}L{3.1cm}L{2.35cm}X@{}}
\toprule
\hrow \thead{From} & \thead{Link} & \thead{To} \\
\midrule
Person & \code{isa} & Mammal \\
Sakti Gauchan & \code{instance} & Person \\
Person & \code{has-part} & Nose \\
Sakti Gauchan & \code{team} & Nepalese team \\
Sakti Gauchan & \code{uniform-color} & Red/Blue \\
\bottomrule
\end{tabularx}
\vspace{2pt}
\begin{itemize}
  \item \textbf{5 sentences, 5 arcs.} The count checks.
  \item \mk{Note to write beside the drawing}: \emph{has-part} sits on
        \textbf{Person}, not on Sakti Gauchan, so \textbf{every} person inherits a nose.
        That is the inheritance the question is testing.
  \item \yr{78 Po}: identical net, \emph{Sakti Gauchan} replaced by \emph{Sandeep
        Lamichane} in the first four arcs. \mk{For the fifth, draw the
        uniform-colour arc off Sandeep and add one line}: \emph{``the paper names Sakti
        Gauchan here; taken as the same person the other four sentences are about.''}
        Say it and you lose nothing.
\end{itemize}}{%
\figT{c5_semnet_sakti.png}
{\footnotesize\color{sub} \mk{The published answer to this exact
question}, from PS Sir's own slide. Two differences from the paper's wording, both
harmless: it draws \code{has\_part} $\rightarrow$ \emph{Head} where the paper says
\emph{nose}, and \code{team\_colours} $\rightarrow$ \emph{Black/Blue} where the paper says
\emph{Red/Blue}. Substitute the paper's words; the shape is unchanged. Source: PS Sir
ch5.}
\lead{The FOPL the deck writes underneath}
\begin{itemize}
  \item \code{ISA(Mammal, Person)}, \code{INSTANCE(Person, Sakti Gauchan)},
        \code{HAS\_PART(Person, head)}, \code{TEAM(Nepal, Sakti Gauchan)},
        \code{TEAM\_COLOURS(black/blue, Sakti Gauchan)}.
  \item \mk{The deck is inconsistent, twice over.} An earlier slide
        writes \code{MAMMAL(person)} for the first sentence $-$ a \emph{unary} predicate
        that loses the relation entirely. And it writes the arguments
        \code{(class, individual)} here, the reverse of Rich and Knight's
        \code{instance(Pee-Wee-Reese, Person)}.
  \item \textbf{Use one order and say which.} This file writes
        \code{isa(subclass, superclass)} and \code{instance(individual, class)}, matching
        the arrow direction in every drawing. \mk{Consistency is what is
        marked, not the choice}: pick one, state it in a line, and never swap.
\end{itemize}}

%=======================================================================
\T{5.N.2 The baseball player: Pee-wee-Reese, and the same set as Rohit Paudel}

\creamq{Convert the given sentences into a semantic net}{\m{7}\m{2+6}}
{\color{sub}\footnotesize \tP{2} \yr{\textbf{73 Bh} \m{7}, \textbf{79 Ch} \m{2+6}}
\gap$\cdot$\gap again \mk{two papers, one net}: \yr{\textbf{79 Ch}}
is \yr{\textbf{73 Bh}} moved from baseball to cricket \gap$\cdot$\gap this is Rich and
Knight's own figure, so \mk{it is also the set to expect in a
``convert the net into a frame'' question}}

\begin{asked}{\textbf{73 Bh} \textperiodcentered{} Q6 \textperiodcentered{} \m{7}}
Convert given sentences into Semantic Network.
\begin{itemize}[topsep=1pt]
  \item i) The height of the adult male is 5.10
  \item ii) Baseball player is an adult male.
  \item iii) Adult male is a person.
  \item iv) Batting average of Baseball players is 0.252
  \item v) Pee-wee-Reese is a Fielder.
  \item vi) Fielder is Baseball player.
  \item vii) Team of pee-wee-Reese is Brooklyn Dodger.
\end{itemize}
\end{asked}

\begin{asked}{\textbf{79 Ch} \textperiodcentered{} Q4 \textperiodcentered{} \m{2+6}}
What is semantic network? Convert the following sentences into semantic network.
\begin{itemize}[topsep=1pt]
  \item i) The avg height of an adult is 5'6''.
  \item ii) Cricket player is an adult male.
  \item iii) Adult male is a person.
  \item iv) Batting average of a cricket player is 37.59
  \item v) Rohit Paudel is a Fielder.
  \item vi) Fielder is a cricket player.
  \item vii) Team of Rohit Paudel is Nepal.
\end{itemize}
{\color{sub} Sentence (i) says \emph{adult} where (ii) and (iii) say \emph{adult male}.
Read as the same node.}
\end{asked}

\sbsr{0.52}{%
\lead{The spine is four deep, and that is the whole difficulty}
\begin{tabularx}{\linewidth}{@{}L{3.0cm}L{2.6cm}X@{}}
\toprule
\hrow \thead{From} & \thead{Link} & \thead{To} \\
\midrule
Adult Male & \code{isa} & Person \\
Baseball Player & \code{isa} & Adult Male \\
Fielder & \code{isa} & Baseball Player \\
Pee-wee-Reese & \code{instance} & Fielder \\
Adult Male & \code{height} & 5.10 \\
Baseball Player & \code{batting-average} & 0.252 \\
Pee-wee-Reese & \code{team} & Brooklyn Dodger \\
\bottomrule
\end{tabularx}
\vspace{2pt}
\begin{itemize}
  \item \textbf{7 sentences, 7 arcs.} The count checks.
  \item \mk{The mark is in where the properties hang.}
        \emph{height} goes on \textbf{Adult Male} and \emph{batting-average} on
        \textbf{Baseball Player}, because that is the class each sentence names. Hanging
        either off Pee-wee-Reese loses the inheritance the question is set to test.
  \item \textbf{Read the inheritance back} in one line: Pee-wee-Reese is a Fielder, so a
        Baseball Player, so an Adult Male, so a Person $-$ and therefore
        \mk{inherits height 5.10 and batting average 0.252 without either
        being stored on him.}
  \item \yr{\textbf{79 Ch}}: same seven arcs. Baseball Player $\rightarrow$ Cricket
        Player, 5.10 $\rightarrow$ 5'6'', 0.252 $\rightarrow$ 37.59, Pee-wee-Reese
        $\rightarrow$ Rohit Paudel, Brooklyn Dodger $\rightarrow$ Nepal.
\end{itemize}}{%
\figT{c5_semnet_peewee.png}
{\footnotesize\color{sub} The figure both papers are drawn from, in its shorter textbook
form: it keeps \emph{Person}, \emph{Pee-Wee-Reese} and the \emph{team} arc but collapses
the Adult Male / Baseball Player / Fielder spine, and its property arcs are
\code{has-part} and \code{uniform-color}. \mk{Copy the shape, use the
seven arcs in the table} for the exam answer. Source: PS Sir ch5.}
\lead{Why this set is worth extra attention}
\begin{itemize}
  \item It is the only sentence set in the papers with \textbf{four levels of class}, so
        it is the one that shows inheritance properly.
  \item It is also \mk{the set \yr{\texttt{81 Ba}} converts into a frame
        system} $-$ see \S5.N.10, where the finished frames are printed.
\end{itemize}}

%=======================================================================
\T{5.N.3 Tom the cat: two sentence sets, three papers}

\creamq{Represent the following sentences using semantic net}{\m{2+5}\m{2+6}}
{\color{sub}\footnotesize \tP{3} \yr{\textbf{\texttt{82 Bh}} \m{2+5},
\textbf{81 Ch} \m{2+6}, \textbf{\texttt{79 Bh}} \m{2+6}} \gap$\cdot$\gap
\mk{the most-repeated set in the chapter} \gap$\cdot$\gap
\yr{\textbf{\texttt{82 Bh}}} and \yr{\textbf{81 Ch}} print it word for word identically;
\yr{\textbf{\texttt{79 Bh}}} is a rewritten variant \gap$\cdot$\gap all three pair it with
a KR short answer, \emph{issues in KR} twice and \emph{define semantic net} once}

\begin{asked}{\textbf{\texttt{82 Bh}} Q5 \m{2+5} \textperiodcentered{} \textbf{81 Ch} Q5 \m{2+6} \textperiodcentered{} the same eleven sentences}
What are different issues in knowledge representation? Represent the following sentences
using semantic net representation.

Tom is a cat. Tom caught a bird. Tom is owned by Sudeep. Tom is ginger in color. Cats like
cream. The cat sat on the mat. A cat is a mammal. A bird is an animal. All mammals are
animals. Mammals have fur. Sudeep age is 20.
\end{asked}

\sbsr{0.46}{%
\figT{c5_semnet_tom.png}
{\footnotesize\color{sub} \mk{The published answer, ten of the eleven
arcs already drawn.} It is the textbook set, so the only difference from the paper is the
owner's name: \emph{John} in the figure, \emph{Sudeep} in the paper. The eleventh
sentence, \emph{Sudeep age is 20}, is not in the figure $-$ add one arc,
\emph{Sudeep} \code{age} $\rightarrow$ \emph{20}. Source: BA Sir CH-05.}}{%
\lead{The eleven arcs, in the order the sentences give them}
\begin{tabularx}{\linewidth}{@{}L{2.1cm}L{2.35cm}X@{}}
\toprule
\hrow \thead{From} & \thead{Link} & \thead{To} \\
\midrule
Tom & \code{instance} & Cat \\
Tom & \code{caught} & Bird \\
Tom & \code{is-owned-by} & Sudeep \\
Tom & \code{is-coloured} & Ginger \\
Cat & \code{like} & Cream \\
Cat & \code{sat-on} & Mat \\
Cat & \code{isa} & Mammal \\
Bird & \code{isa} & Animal \\
Mammal & \code{isa} & Animal \\
Mammal & \code{has} & Fur \\
Sudeep & \code{age} & 20 \\
\bottomrule
\end{tabularx}
\vspace{2pt}
\begin{itemize}
  \item \textbf{11 sentences, 11 arcs.} The count checks.
  \item \mk{Two traps, and both are worth a mark each}:
  \begin{itemize}
    \item \emph{Tom is a cat} is \code{instance} (Tom is an individual);
          \emph{a cat is a mammal} is \code{isa} (both are classes). Same English
          construction, different link.
    \item \emph{The cat sat on the mat} is about \textbf{the class} Cat as the figure draws
          it, but the sentence plainly means \emph{that} cat, ie Tom.
          \mk{Draw it as the figure does and add a line saying the
          natural reading is Tom} $-$ the ambiguity is exactly the ``lack of formality''
          limitation the paired theory question asks about.
  \end{itemize}
  \item \emph{Cats like cream} and \emph{mammals have fur} are drawn out of the
        \textbf{class} nodes, so Tom inherits both. Say so.
\end{itemize}}

\creamq{The 2079 Bhadra variant: eleven different sentences, same shape}{\m{2+6}}
{\color{sub}\footnotesize \yr{\textbf{\texttt{79 Bh}}} \m{2+6} \gap$\cdot$\gap counted in
the chip above \gap$\cdot$\gap \mk{the last sentence is the interesting
one}: it is a quantified statement, which a semantic net cannot properly express}

\begin{asked}{\textbf{\texttt{79 Bh}} \textperiodcentered{} Q6 \textperiodcentered{} \m{2+6}}
List the issues need to be consider in Knowledge Representation techniques. Convert given
sentences into Semantic Network.
\begin{itemize}[topsep=1pt]
  \item a) Tom is a cat \quad b) Tom fights with rat \quad c) Tom is owned by Ram
  \item d) Tom is black in color. \quad e) Cats like milk \quad f) Rats like cheese
  \item g) The cat sat on the bed \quad h) A cat is a mammal \quad i) A Rat is an animal
  \item j) All mammals are animals \quad k) Every mammal gives birth a baby
\end{itemize}
\end{asked}

\sbs{%
\begin{tabularx}{\linewidth}{@{}L{2.0cm}L{2.6cm}X@{}}
\toprule
\hrow \thead{From} & \thead{Link} & \thead{To} \\
\midrule
Tom & \code{instance} & Cat \\
Tom & \code{fights-with} & Rat \\
Tom & \code{is-owned-by} & Ram \\
Tom & \code{is-coloured} & Black \\
Cat & \code{like} & Milk \\
Rat & \code{like} & Cheese \\
Cat & \code{sat-on} & Bed \\
Cat & \code{isa} & Mammal \\
Rat & \code{isa} & Animal \\
Mammal & \code{isa} & Animal \\
Mammal & \code{gives-birth-to} & Baby \\
\bottomrule
\end{tabularx}}{%
\begin{itemize}
  \item \textbf{11 sentences, 11 arcs.} The count checks.
  \item Layout: \emph{Animal} at the top, \emph{Mammal} and \emph{Rat} under it,
        \emph{Cat} under \emph{Mammal}, \emph{Tom} at the bottom, and every value node
        (Milk, Cheese, Bed, Black, Ram, Baby) hanging sideways.
  \item \mk{(k) is the one that earns the extra mark.}
        \emph{``Every mammal gives birth to a baby''} is
        $\forall x\,(x \in Mammal \Rightarrow \exists y\,(y \in Baby \wedge
        gives\text{-}birth\text{-}to(x,y)))$.
  \begin{itemize}
    \item The net can only draw \textbf{one arc} out of the \emph{Mammal} class node.
    \item \textbf{Write underneath}: \emph{``the arc out of a class node is read as
          `for every member', and the existential `some baby' is not expressible in a
          semantic net at all $-$ this is the quantification limitation.''}
  \end{itemize}
  \item \emph{Tom fights with rat} is between Tom and the \textbf{class} Rat as printed.
        Note that a fight is with one particular rat, so a strict reading needs an
        \textbf{event node} (\S5.4 of the theory).
\end{itemize}}

%=======================================================================
\T{5.N.4 Mammals, fur and water}

\creamq{Define term frame. Convert the following sentences into a semantic network}{\m{1+6}}
{\color{sub}\footnotesize \tP{1} \yr{82 Ka \m{1+6}} \gap$\cdot$\gap
\mk{the set with three siblings under one class}, so it is the one that
rewards pushing a property up \gap$\cdot$\gap the 1-mark opener is
\emph{define a frame}, \S5.5 of the theory, one sentence}

\begin{asked}{82 Ka \textperiodcentered{} Q5 \textperiodcentered{} \m{1+6}}
Define term frame. Convert the following sentences into semantic network. Cat is a mammal.
Cat has fur. Mammal has vertebra. Mammal is an animal. Bear is a mammal. Bear has fur.
Whale is a mammal. Whale lives in water. \mk{Fish is an animal 10.} Fish
lives in water.

{\color{sub} The \emph{``10.''} is the paper's own stray list number printed inside the
sentence. Reproduced as printed; it is not a tenth sentence.}
\end{asked}

\sbs{%
\begin{tabularx}{\linewidth}{@{}L{1.9cm}L{2.0cm}X@{}}
\toprule
\hrow \thead{From} & \thead{Link} & \thead{To} \\
\midrule
Cat & \code{isa} & Mammal \\
Bear & \code{isa} & Mammal \\
Whale & \code{isa} & Mammal \\
Mammal & \code{isa} & Animal \\
Fish & \code{isa} & Animal \\
Cat & \code{has} & Fur \\
Bear & \code{has} & Fur \\
Mammal & \code{has} & Vertebra \\
Whale & \code{lives-in} & Water \\
Fish & \code{lives-in} & Water \\
\bottomrule
\end{tabularx}
\vspace{2pt}
\begin{itemize}
  \item \textbf{10 sentences, 10 arcs.} The count checks.
  \item \mk{Every node here is a class}, so \textbf{every} link is
        \code{isa} or a property arc. There is not one \code{instance} arc in this set $-$
        no individual is named. Saying that in one line shows you understood the
        difference rather than guessed it.
\end{itemize}}{%
\lead{Layout}
\begin{itemize}
  \item \emph{Animal} on the top row. \emph{Mammal} and \emph{Fish} on the second.
        \emph{Cat}, \emph{Bear}, \emph{Whale} on the third. Values (\emph{Fur},
        \emph{Vertebra}, \emph{Water}) off to the right.
  \item \emph{Water} is one node reached by two arcs, not two nodes.
        \mk{A semantic net stores each object once}; drawing it twice is
        the mistake that costs a mark here.
\end{itemize}
\lead{\mk{The observation that earns the last mark}}
\begin{itemize}
  \item \emph{Cat has fur} and \emph{Bear has fur} are both given, and both are drawn
        because the question says to convert the sentences.
  \item But fur is a \textbf{mammal} property. Add a line: \emph{``\code{has fur} could be
        moved onto Mammal, and then Cat, Bear and Whale would all inherit it $-$ that is
        the space saving inheritance buys.''}
  \item \mk{Then give the counter-example in the same set}: whales are
        mammals and do \emph{not} have fur in the ordinary sense, so moving it up would
        make the net assert something false. \textbf{A net has no way to state an
        exception}, which is precisely why frames carry defaults and overrides.
  \item That pair of sentences is the best two-line answer to \emph{``limitations of
        semantic nets''} anywhere in these papers, and it comes free with this question.
\end{itemize}}

%=======================================================================
\T{5.N.5 Fido the dog}

\creamq{Explain the approaches of KR. Highlight pros and cons of semantic nets and frames. Use them to represent the following sentences}{\m{2+6}}
{\color{sub}\footnotesize \tP{1} \yr{\textbf{\texttt{81 Bh}} \m{2+6}} \gap$\cdot$\gap
the smallest set in the papers, \textbf{six sentences} \gap$\cdot$\gap
\mk{it asks for both representations}, so draw the net \emph{and} write
the frames $-$ the arc list below does both}

\begin{asked}{\textbf{\texttt{81 Bh}} \textperiodcentered{} Q6 \textperiodcentered{} \m{2+6}}
Explain different approaches of knowledge representation. Highlight pros and cons of
Semantic Nets and Frames. Use them to represent following sentences.
\begin{itemize}[topsep=1pt]
  \item \mk{Figo is a dog.}
  \item Fido is own by Rohit.
  \item Fido is ginger in color.
  \item Dog is a mammal.
  \item All Mammals are animals.
  \item Mammals have fur.
\end{itemize}
{\color{sub} The paper names the dog \emph{Figo} in the first line and \emph{Fido} in the
next two. Reproduced as printed; take them as one animal and say so.}
\end{asked}

\sbsr{0.47}{%
\lead{As a semantic net}
\begin{tabularx}{\linewidth}{@{}L{1.7cm}L{2.35cm}X@{}}
\toprule
\hrow \thead{From} & \thead{Link} & \thead{To} \\
\midrule
Fido & \code{instance} & Dog \\
Fido & \code{is-owned-by} & Rohit \\
Fido & \code{is-coloured} & Ginger \\
Dog & \code{isa} & Mammal \\
Mammal & \code{isa} & Animal \\
Mammal & \code{has} & Fur \\
\bottomrule
\end{tabularx}
\vspace{2pt}
\begin{itemize}
  \item \textbf{6 sentences, 6 arcs.} The count checks.
  \item Spine: \emph{Animal} $\leftarrow$ \emph{Mammal} $\leftarrow$ \emph{Dog}, with
        \emph{Fido} hung off \emph{Dog}. \emph{Rohit} and \emph{Ginger} sideways.
  \item \mk{Open the answer with one line on the name}: \emph{``the
        paper writes Figo once and Fido twice; taken as the same dog throughout.''}
        It costs nothing and shows you read the question.
  \item \emph{Mammals have fur} is on the \textbf{class}, so \mk{Fido
        inherits fur} without it being stored on him. Write that sentence out $-$ this
        set is small enough that the inheritance remark is where the marks are.
\end{itemize}}{%
\lead{The same knowledge as frames}
\begin{tabularx}{\linewidth}{@{}L{1.75cm}L{2.0cm}X@{}}
\toprule
\hrow \thead{Frame} & \thead{Slot} & \thead{Filler} \\
\midrule
\textbf{Animal} & \emph{(class frame)} & $-$ \\
\midrule
\textbf{Mammal} & \code{isa} & Animal \\
 & \code{has} & Fur \\
\midrule
\textbf{Dog} & \code{isa} & Mammal \\
\midrule
\textbf{Fido} & \code{instance} & Dog \\
 & \code{owner} & Rohit \\
 & \code{colour} & Ginger \\
 & \code{has} & \emph{Fur} (inherited) \\
\bottomrule
\end{tabularx}
\vspace{2pt}
\begin{itemize}
  \item \textbf{One frame per node, one slot per outgoing arc}, exactly the recipe in
        \S5.6 of the theory.
  \item \mk{Write the inherited row in italics or brackets} and label it
        \emph{inherited}. It shows the frame system doing the work, and it is the
        difference between copying the net out as a table and answering the question.
  \item \textbf{Pros and cons} (the first half of the question) are the two lists in
        \S5.4 and \S5.5 of the theory. One line each is enough at 2 marks.
\end{itemize}}

%=======================================================================
\T{5.N.6 Employees and supervisors}

\creamq{When do we use a semantic network? Convert the following sentences into a semantic network}{\m{2+6}}
{\color{sub}\footnotesize \tP{1} \yr{\texttt{82 Ba} \m{2+6}} \gap$\cdot$\gap
\mk{the trickiest set in the papers}: it contains a word that is both a
class and a relation, and a symmetric relation \gap$\cdot$\gap the 2-mark opener,
\emph{when do we use a semantic net}, is \S5.4 of the theory}

\begin{asked}{\texttt{82 Ba} \textperiodcentered{} Q6 \textperiodcentered{} \m{2+6}}
When do we use semantic network? Convert the following sentences into semantic network.
\begin{itemize}[topsep=1pt]
  \item i) Employee is a human. \quad ii) Hari is an employee. \quad
        iii) Sita is a human.
  \item iv) Supervisor is a human. \quad v) Ram is a supervisor. \quad
        vi) Ram is married to Sita.
  \item vii) Ram is supervisor of Hari.
\end{itemize}
\end{asked}

\sbs{%
\begin{tabularx}{\linewidth}{@{}L{1.6cm}L{2.5cm}X@{}}
\toprule
\hrow \thead{From} & \thead{Link} & \thead{To} \\
\midrule
Employee & \code{isa} & Human \\
Supervisor & \code{isa} & Human \\
Hari & \code{instance} & Employee \\
Sita & \code{instance} & Human \\
Ram & \code{instance} & Supervisor \\
Ram & \code{married-to} & Sita \\
Ram & \code{supervisor-of} & Hari \\
\bottomrule
\end{tabularx}
\vspace{2pt}
\begin{itemize}
  \item \textbf{7 sentences, 7 arcs.} The count checks.
  \item Layout: \emph{Human} at the top; \emph{Employee} and \emph{Supervisor} under it;
        \emph{Hari}, \emph{Ram}, \emph{Sita} on the bottom row so the two arcs between
        individuals are horizontal and easy to read.
\end{itemize}}{%
\lead{\mk{Three things the examiner is testing here}}
\begin{itemize}
  \item \textbf{1. \emph{Supervisor} is used two ways in one question.}
  \begin{itemize}
    \item In (iv) and (v) it is a \textbf{class} node, reached by \code{isa} and
          \code{instance}.
    \item In (vii) \code{supervisor-of} is a \textbf{relation arc between two
          individuals}, Ram to Hari. \mk{It is not the class node.}
    \item Drawing an arc from Hari to the \emph{Supervisor} class is the mistake this
          set is built to catch.
  \end{itemize}
  \item \textbf{2. Sita is an instance of \emph{Human} directly}, not of Employee or
        Supervisor. The sentences never say what she does, so do not invent a class for
        her. \mk{Adding information the sentences do not give loses
        marks}, however sensible it looks.
  \item \textbf{3. \code{married-to} is symmetric.} A semantic net arc is directed, so
        either draw a \textbf{double-headed} arrow or add the inverse arc
        \emph{Sita} \code{married-to} $\rightarrow$ \emph{Ram}, and say which convention
        you used.
  \begin{itemize}
    \item This is the \emph{``relationships among attributes''} issue from \S5.3: inverse
          and symmetric relations have to be handled explicitly, because the formalism
          does not know they exist.
  \end{itemize}
\end{itemize}}

%=======================================================================
\T{5.N.7 FOPL first, then the net}

\creamq{Convert the following sentences into FOPL and represent using a semantic network}{\m{2+6}}
{\color{sub}\footnotesize \tP{1} \yr{81 Ash \m{2+6}} \gap$\cdot$\gap
\mk{the only set that asks for FOPL as well as the drawing}, so it is
worth doing every set this way in practice \gap$\cdot$\gap the FOPL is free marks: it is
the arc list written as predicates}

\begin{asked}{81 Ash \textperiodcentered{} Q6 \textperiodcentered{} \m{2+6}}
What are Frames and Sematic Network? Convert the following sentences into FOPL and
represent using semantic network.
\begin{itemize}[topsep=1pt]
  \item i) Ram is a person. \quad ii) Person is an animal. \quad
        iii) Ram likes Jhilke. \quad iv) Jhilke is a cat.
  \item v) Cat eats fish. \quad vi) Cat is an animal. \quad vii) fish is an animal.
\end{itemize}
\end{asked}

\sbsr{0.53}{%
\lead{Arc list and FOPL, side by side}
\begin{tabularx}{\linewidth}{@{}L{1.35cm}L{1.85cm}L{1.35cm}X@{}}
\toprule
\hrow \thead{From} & \thead{Link} & \thead{To} & \thead{FOPL} \\
\midrule
Ram & \code{instance} & Person & \code{instance(Ram, Person)} \\
Person & \code{isa} & Animal & \code{isa(Person, Animal)} \\
Ram & \code{likes} & Jhilke & \code{likes(Ram, Jhilke)} \\
Jhilke & \code{instance} & Cat & \code{instance(Jhilke, Cat)} \\
Cat & \code{eats} & Fish & \emph{see below} \\
Cat & \code{isa} & Animal & \code{isa(Cat, Animal)} \\
Fish & \code{isa} & Animal & \code{isa(Fish, Animal)} \\
\bottomrule
\end{tabularx}
\vspace{2pt}
\begin{itemize}
  \item \textbf{7 sentences, 7 arcs.} The count checks.
  \item \mk{The rule for turning an arc into FOPL}: an arc
        $A \xrightarrow{R} B$ becomes the binary predicate $R(A,B)$, with \code{isa} used
        between classes and \code{instance} between an individual and its class. Nothing
        else is needed for six of the seven.
\end{itemize}}{%
\lead{\mk{Sentence (v) is the one worth thinking about}}
\begin{itemize}
  \item \emph{``Cat eats fish''} is a statement about \textbf{every} cat and
        \textbf{some} fish, so its honest FOPL is
        $$\forall x\,\big(Cat(x) \Rightarrow \exists y\,(Fish(y) \wedge eats(x,y))\big)$$
  \item The net can only draw \textbf{one arc}, \emph{Cat} \code{eats}
        $\rightarrow$ \emph{Fish}, class to class.
  \item \mk{Write both, and say what is lost}: the arc records the
        relation but throws away the quantifiers. That single remark is the entire
        \emph{``a semantic net cannot express quantification''} limitation, evidenced on
        the examiner's own data.
  \item If you want the flat version for the arc table, \code{eats(Cat, Fish)} is what
        every deck writes and is accepted.
\end{itemize}
\lead{Layout}
\begin{itemize}
  \item \emph{Animal} on top, with \textbf{three} \code{isa} arcs into it, from
        \emph{Person}, \emph{Cat} and \emph{Fish}.
  \item \emph{Ram} under \emph{Person}, \emph{Jhilke} under \emph{Cat}, and the
        \code{likes} arc running horizontally between the two individuals.
  \item \mk{The net has a nice shape here}: two parallel
        individual-to-class chains joined at the top by \emph{Animal} and at the bottom by
        \code{likes}. Draw it that way and the seven arcs never cross.
\end{itemize}}

%=======================================================================
\T{5.N.8 Birds: Tweety and Sweety}

\creamq{What are frames and semantic nets? Convert the given sentences into a semantic net}{\m{2+6}}
{\color{sub}\footnotesize \tP{1} \yr{\textbf{\textit{77 Ch}} \m{2+6}} \gap$\cdot$\gap
\mk{two sentences here need more than one arc each}, which is what
makes this set harder than it looks \gap$\cdot$\gap the same paper asks conceptual
dependency and scripts at Q\,8, so \S5.7 of the theory is on the same sheet}

\begin{asked}{\textbf{\textit{77 Ch}} \textperiodcentered{} Q6 \textperiodcentered{} \m{2+6}}
What are Frames and Semantic Net? Convert the given sentences in semantic Net:
\begin{itemize}[topsep=1pt]
  \item Tweety and Sweety are birds. \quad Tweety has a red beak. \quad
        Sweety is Tweety's child.
  \item A crow is a bird. \quad Birds can fly. \quad Sparrow is a bird. \quad
        Sparrow has a wing.
\end{itemize}
\end{asked}

\sbsr{0.55}{%
\begin{tabularx}{\linewidth}{@{}L{1.6cm}L{2.2cm}X@{}}
\toprule
\hrow \thead{From} & \thead{Link} & \thead{To} \\
\midrule
Tweety & \code{instance} & Bird \\
Sweety & \code{instance} & Bird \\
Tweety & \code{has-part} & Beak \\
Beak & \code{is-coloured} & Red \\
Sweety & \code{child-of} & Tweety \\
Tweety & \code{parent-of} & Sweety \quad {\footnotesize\color{sub}(inverse)} \\
Crow & \code{isa} & Bird \\
Sparrow & \code{isa} & Bird \\
Bird & \code{can} & Fly \\
Sparrow & \code{has-part} & Wing \\
\bottomrule
\end{tabularx}
\vspace{2pt}
\begin{itemize}
  \item \textbf{7 sentences, 10 arcs.} Here the count does \emph{not} match one to one,
        and that is the point:
  \begin{itemize}
    \item \emph{Tweety and Sweety are birds} is \textbf{two} \code{instance} arcs, one
          per bird.
    \item \emph{Tweety has a red beak} is \textbf{two} arcs: the beak is a part of Tweety,
          and \emph{red} is a property of the \textbf{beak}, not of Tweety.
          \mk{Drawing one arc \emph{Tweety} \code{red-beak}
          $\rightarrow$ \emph{Red} loses the mark} $-$ the colour belongs to the part.
    \item \emph{Sweety is Tweety's child} is one arc, but the \textbf{inverse} arc is
          worth drawing and labelling.
  \end{itemize}
\end{itemize}}{%
\figT{c5_semnet_bird.png}
{\footnotesize\color{sub} \mk{The shape to copy}, from PS Sir's deck:
\emph{Animal} $\leftarrow$ \emph{Bird} $\leftarrow$ \emph{Robin} $\leftarrow$
\emph{Rusty} / \emph{Red}, with \code{hasPart} $\rightarrow$ \emph{Wing} drawn once on
the \emph{Bird} class. Substitute this paper's nodes. Source: PS Sir ch5.}
\begin{itemize}
  \item \mk{Note to add}: \emph{Sparrow has a wing} is drawn as the
        paper gives it, on \emph{Sparrow}, but wings are a property of \textbf{all} birds
        $-$ moving \code{has-part} $\rightarrow$ \emph{Wing} up to the \emph{Bird} class
        would let Tweety, Sweety, Crow and Sparrow inherit it from one arc.
  \item \emph{Birds can fly} is on the class, so every bird inherits it. That is also the
        standard \textbf{exception} example: a penguin is a bird and cannot fly, and a
        semantic net has no way to override an inherited property. Frames do, with a
        default.
\end{itemize}}

%=======================================================================
\T{5.N.9 A semantic net about the AI course itself}

\creamq{Draw a sample net about your course AI to demonstrate the features and qualities of a semantic net}{\m{7}}
{\color{sub}\footnotesize \tP{1} \yr{\textit{72 Ma} \m{7}} \gap$\cdot$\gap
\mk{the only open-ended drawing question in the papers}: the content is
yours, the marks are for demonstrating the \emph{features} \gap$\cdot$\gap so build the
net so that each feature appears at least once, and label them}

\begin{asked}{\textit{72 Ma} \textperiodcentered{} Q5 \textperiodcentered{} \m{7}}
What is a semantic net? Draw a sample net about your course AI to demonstrate the
knowledge representation features and qualities of semantic net.
\end{asked}

\sbs{%
\lead{A net that hits every feature}
\begin{tabularx}{\linewidth}{@{}L{2.5cm}L{2.35cm}X@{}}
\toprule
\hrow \thead{From} & \thead{Link} & \thead{To} \\
\midrule
Elective & \code{isa} & Course \\
AI (CT 710) & \code{instance} & Course \\
Course & \code{has-part} & Chapter \\
AI (CT 710) & \code{credit-hours} & 3 \\
AI (CT 710) & \code{semester} & IV/I \\
AI (CT 710) & \code{programme} & BEI \\
Search & \code{part-of} & AI (CT 710) \\
A* & \code{isa} & Search \\
Student & \code{studies} & AI (CT 710) \\
Ram & \code{instance} & Student \\
AI (CT 710) & \code{taught-by} & Teacher \\
Teacher & \code{isa} & Person \\
Student & \code{isa} & Person \\
\bottomrule
\end{tabularx}}{%
\lead{\mk{Label the features on the drawing}, one arrow each}
\begin{itemize}
  \item \textbf{Class hierarchy} $-$ \emph{Person} $\leftarrow$ \emph{Student} /
        \emph{Teacher}, joined by \code{isa}.
  \item \textbf{Instances} $-$ \emph{Ram} \code{instance} $\rightarrow$ \emph{Student}.
        Point out that \code{isa} and \code{instance} are different links.
  \item \textbf{Inheritance} $-$ anything put on \emph{Person} is available to
        \emph{Ram} through two arcs. Write the chain out.
  \item \textbf{Part-whole} $-$ \emph{Search} \code{part-of} $\rightarrow$ \emph{AI},
        which does \textbf{not} inherit the way \code{isa} does. Saying that is worth a
        mark on its own.
  \item \textbf{Attribute arcs} $-$ \emph{credit-hours}, \emph{semester},
        \emph{programme}: values rather than objects.
  \item \textbf{Intersection search} $-$ start from \emph{Ram} and from \emph{A*} and the
        two wavefronts meet at \emph{AI (CT 710)}, which answers
        \emph{``what connects Ram and A*?''} without anyone having asked that question in
        advance.
  \item \mk{Close on the limitation}, because the question says
        \emph{qualities}: nothing in this net can say \emph{``no student takes more than
        six electives''} or \emph{``AI is not a lab course''}. No quantifiers, no
        negation.
\end{itemize}}

%=======================================================================
\T{5.N.10 Converting a semantic net into a frame}

\creamq{Compare semantic net and frame. How do you convert a semantic net into a frame? Explain with example}{\m{3+5}}
{\color{sub}\footnotesize \tP{1} \yr{\texttt{81 Ba} \m{3+5}} \gap$\cdot$\gap
the comparison table is \S5.6 of the theory \gap$\cdot$\gap
\mk{worked here on the Pee-wee-Reese net of \S5.N.2}, because that is
the set the textbook itself converts, and both halves are published figures}

\begin{asked}{\texttt{81 Ba} \textperiodcentered{} Q6 \textperiodcentered{} \m{3+5}}
Compare semantic net and frame. How you convert semantic net into frame? Explain with
example.
\end{asked}

\sbsr{0.44}{%
\lead{The five steps, applied}
\begin{itemize}
  \item \textbf{1. Every node becomes a frame.} The net of \S5.N.2 has seven nodes, so the
        frame system has seven frames: \emph{Person}, \emph{Adult-Male},
        \emph{ML-Baseball-Player}, \emph{Fielder}, \emph{Pee-Wee-Reese},
        \emph{ML-Baseball-Team}, \emph{Brooklyn-Dodgers}.
  \item \textbf{2. Every outgoing arc becomes a slot.} Arc label $\rightarrow$ slot name,
        arc target $\rightarrow$ filler. \emph{Adult-Male} \code{height} $\rightarrow$
        \emph{5-10} becomes a \code{height} slot in the \emph{Adult-Male} frame with
        filler \emph{5-10}.
  \item \textbf{3. \code{isa} and \code{instance} become the class link}, not ordinary
        slots. \emph{Fielder} gets \code{isa: ML-Baseball-Player} and is a
        \textbf{subclass} frame; \emph{Pee-Wee-Reese} gets \code{instance: Fielder} and is
        an \textbf{instance} frame.
  \item \textbf{4. A filler that has its own frame stays a pointer.} \emph{Pee-Wee-Reese}
        \code{team} $\rightarrow$ \emph{Brooklyn-Dodgers} is a link to that frame, not the
        string. \mk{This is what makes the result a frame system rather
        than seven separate tables.}
  \item \textbf{5. Push each property to the highest class that has it}, and let everything
        below inherit. \code{height} sits on \emph{Adult-Male},
        \code{batting-average} on \emph{ML-Baseball-Player}, and Pee-Wee-Reese overrides
        both with his own values.
\end{itemize}
\lead{\mk{Then check, and say you checked}}
\begin{itemize}
  \item \textbf{No information may be lost.} Count the arcs in the net, count the filled
        slots in the frames. They must agree.
  \item \textbf{What the frames gained}: \emph{cardinality} (6,000,000,000 persons),
        \emph{defaults} (\emph{handed: Right}), and slots marked with \code{*} for values
        that belong to the class rather than to each member. None of those can be drawn in
        a net.
\end{itemize}}{%
\figT{c5_frame_system.png}
{\footnotesize\color{sub} \mk{The net of \S5.N.2 after conversion.}
Seven frames, each slot either an inherited class link, a value, or a pointer to another
frame. The \code{*} marks a value stored on the class. Compare with
\code{c5\_semnet\_peewee} above: same knowledge, different formalism. Source: PS Sir ch5,
after Rich and Knight Fig 9.5.}
\lead{The reverse direction, if asked}
\begin{itemize}
  \item Frame $\rightarrow$ net is step 1 to 4 run backwards: one node per frame, one arc
        per filled slot, and \code{isa} / \code{instance} arcs from the class links.
  \item \mk{The conversion is lossy in this direction.} Defaults,
        cardinalities, type restrictions and demons have no arc to become, so they are
        simply dropped. Say so $-$ it is the cleanest one-line proof that a frame carries
        more than a net.
\end{itemize}}

%=======================================================================
\T{5.N.11 Semantic net and frame with FOPL statements}

\creamq{Why are semantic networks and frames important in AI? Provide examples of both with FOPL statements}{\m{2+6}}
{\color{sub}\footnotesize \tP{1} \yr{75 Ba \m{2+6}} \gap$\cdot$\gap
\mk{three formalisms in one answer} $-$ net, frame and FOPL, all saying
the same thing \gap$\cdot$\gap the 2-mark opener is the \emph{why they matter} list in
\S5.6 of the theory}

\begin{asked}{75 Ba \textperiodcentered{} Q6 \textperiodcentered{} \m{2+6}}
Why semantic network and frames are important in AI? Provide examples of both with FOPL
statements example.
\end{asked}

\sbsr{0.45}{%
\figT{c5_frame_fopl.png}
{\footnotesize\color{sub} (a) the frame-based knowledge base and (b) the first-order
sentences that say the same thing. The semantic net is the same picture with the boxes
redrawn as nodes and the \code{Subset} / \code{Member} labels put on the arcs. Source:
PS Sir ch5.}}{%
\lead{The three-way translation table to reproduce}
\begin{tabularx}{\linewidth}{@{}L{2.15cm}L{2.5cm}X@{}}
\toprule
\hrow \thead{Net arc} & \thead{Frame} & \thead{FOPL} \\
\midrule
$A \xrightarrow{isa} B$ & \code{isa: B} in frame $A$ & \code{Subset(A, B)}, ie
$A \subset B$ \\
$A \xrightarrow{instance} B$ & \code{instance: B} in frame $A$ & \code{Member(A, B)}, ie
$A \in B$ \\
$A \xrightarrow{R} v$ on an individual & slot \code{R} with filler $v$ &
\code{Rel(R, A, v)}, ie $R(A,v)$ \\
$A \xrightarrow{R} v$ on a class & slot \code{R} on the class frame &
$\forall x\,(x \in A \Rightarrow R(x,v))$ \\
\bottomrule
\end{tabularx}
\vspace{2pt}
\begin{itemize}
  \item Worked on the figure's own data: \code{Subset(Birds, Animals)},
        \code{Subset(Mammals, Animals)}, \code{Member(Opus, Penguins)},
        \code{Rel(Legs, Birds, 2)}, \code{Rel(Legs, Mammals, 4)},
        \code{Rel(Flies, Penguins, F)}, \code{Friend(Opus, Bill)}.
  \item \mk{The answer structure that scores}: draw the small net, write
        the same knowledge as two or three frames beside it, then write one FOPL line per
        arc underneath. Three formalisms, one set of facts, and the marks are for showing
        they are interchangeable.
  \item \mk{Close with what FOPL has that neither picture does}:
        quantifiers, negation and disjunction $-$ and what the pictures have that FOPL
        does not: inheritance for free, defaults, and a shape a human can check by eye.
        That is the \emph{``why they are important in AI''} half of the question.
\end{itemize}}

\hr

{\color{sub}\footnotesize \textbf{Coverage.} All \textbf{12 papers} that set a
sentence-set conversion are worked above: \yr{72 Ma}, \yr{78 Po} (\S5.N.1);
\yr{\textbf{73 Bh}}, \yr{\textbf{79 Ch}} (\S5.N.2); \yr{\textbf{\texttt{82 Bh}}},
\yr{\textbf{81 Ch}}, \yr{\textbf{\texttt{79 Bh}}} (\S5.N.3); \yr{82 Ka} (\S5.N.4);
\yr{\textbf{\texttt{81 Bh}}} (\S5.N.5); \yr{\texttt{82 Ba}} (\S5.N.6); \yr{81 Ash}
(\S5.N.7); \yr{\textbf{\textit{77 Ch}}} (\S5.N.8). \S5.N.9 to \S5.N.11 work the three
remaining drawing questions: the open-ended net \yr{\textit{72 Ma}}, the net-to-frame
conversion \yr{\texttt{81 Ba}}, and the net/frame/FOPL question \yr{75 Ba}.
\gap$\cdot$\gap \textbf{Three papers print a defect in their own sentence list}, all
reproduced above as printed and all flagged: \yr{78 Po} names Sakti Gauchan in a Sandeep
Lamichane set, \yr{82 Ka} leaves a stray \emph{``10.''} inside a sentence, and
\yr{\textbf{\texttt{81 Bh}}} calls the dog Figo once and Fido twice.
\gap$\cdot$\gap \textbf{Nothing in this file is machine-checkable}: a semantic net has no
arithmetic to recompute, so every arc list above was checked by hand against the OCR
archive, sentence by sentence, and the arc counts are printed with each answer so the
reader can check them too.}
